\section{Анализ предметной области}
\subsection{Характеристика предметной области онлайн-бронирования услуг}

Предметной областью данной работы является организация онлайн-бронирования услуг между частным специалистом (мастером) и клиентом. К таким услугам относятся консультации, персональные занятия, услуги индустрии красоты, обучения и другие форматы, в которых запись производится на фиксированные временные интервалы.

Ключевой особенностью данной предметной области является высокая чувствительность к корректности временных данных. Ошибки в расписании напрямую приводят к конфликтам между клиентами, потере дохода и снижению доверия к исполнителю. По этой причине система записи должна обеспечивать единый источник достоверных данных о доступности мастера и состоянии каждой записи.

В рамках рассматриваемой модели выделяются две основные предметные роли:
\begin{itemize}
\item мастер, который публикует доступные интервалы, управляет записями и взаимодействует с клиентской базой;
\item клиент, который получает доступ к мастеру по приглашению, просматривает доступные слоты и оформляет бронирование.
\end{itemize}

Дополнительно в системе предусмотрена служебная роль администратора. Она не участвует в основном процессе записи, но используется для контроля состояния платформы, просмотра пользователей, записей, аудита и эксплуатационных метрик.

Базовый бизнес-процесс включает следующие этапы: регистрация пользователя, формирование рабочей доступности мастера, предоставление ссылки-приглашения клиенту, выбор клиентом свободного слота, создание бронирования, уведомление сторон, а также возможная отмена или перенос записи.

Для корректной реализации данного процесса требуется формализованная работа с основными сущностями: пользователь, профиль мастера, профиль клиента, приглашение, бронирование, подключение внешнего календаря и события интеграционного обмена. При этом должны соблюдаться базовые ограничения предметной области:
\begin{itemize}
\item недопустимость пересечения активных записей у одного мастера;
\item обязательное выполнение условия начала и окончания интервала записи (\texttt{start\_at < end\_at});
\item контроль прав доступа при изменении и отмене бронирований.
\end{itemize}

\subsection{Проблематика традиционного подхода к организации записи клиентов}

В традиционной организации записи, основанной на сообщениях в мессенджерах, телефонных звонках и ручном ведении расписания, наблюдается ряд системных недостатков.

Во-первых, отсутствует централизованное хранилище информации о клиентах и записях. Данные распределяются по разным каналам связи, что усложняет поиск истории взаимодействия и повышает вероятность ошибок.

Во-вторых, высокий риск двойного бронирования возникает из-за ручного подтверждения интервалов и несвоевременного обновления календаря. При одновременных запросах от нескольких клиентов мастер вынужден самостоятельно контролировать актуальность расписания.

В-третьих, значительная часть рабочего времени тратится на операционные действия: подтверждение записи, согласование переносов, отправку напоминаний и повторные уточнения по времени. Эти операции не создают прямой ценности услуги, но формируют дополнительную нагрузку.

В-четвертых, при отсутствии интеграции с внешним календарем недоступно автоматическое исключение занятых интервалов. В результате клиент может видеть слоты, которые фактически недоступны из-за других встреч мастера.

В-пятых, отсутствие формализованной модели доступа и трассируемости изменений затрудняет анализ инцидентов. При возникновении конфликтов сложно установить, кто и когда изменил конкретную запись.

Таким образом, традиционный подход не обеспечивает необходимый уровень надежности, управляемости и масштабируемости процесса онлайн-записи.

\subsection{Обоснование необходимости автоматизации}

Автоматизация процесса бронирования необходима для перехода от ручного управления записями к цифровой системе, в которой клиентам проще записываться, а мастеру -- управлять расписанием.

Функционально система должна обеспечивать:
\begin{itemize}
\item централизованное хранение данных о пользователях, приглашениях и записях;
\item ролевой доступ для мастера, клиента и администратора;
\item публикацию и расчет доступных слотов с учетом внутренних правил расписания;
\item автоматическое исключение занятых интервалов на основании данных приложения и внешнего календаря;
\item быстрые операции создания, отмены и изменения бронирований;
\item уведомление участников процесса о ключевых событиях жизненного цикла записи.
\end{itemize}

Следовательно, автоматизация процесса онлайн-записи является обоснованным способом решения задач рассматриваемой предметной области.

\subsection{Анализ существующих решений}

В рамках предпроектного обследования предметной области был проведен полевой сбор данных в городе Курске. В ходе обследования посещены парикмахерские и салоны красоты, а также проведены интервью с мастерами о текущей практике ведения записи клиентов. Общая выборка составила порядка 30 мастеров.

По результатам обследования установлено следующее распределение используемых подходов:
\begin{itemize}
\item 15 мастеров (50~\%) ведут запись вручную: бумажный блокнот и подтверждение времени по телефонным звонкам;
\item 8 мастеров (26,7~\%) ведут запись через мессенджеры (Telegram, WhatsApp), без единого календаря и формального контроля конфликтов;
\item 4 мастера (13,3~\%) используют комбинированный подход: бумажный журнал вместе с электронными заметками или таблицами;
\item 3 мастера (10~\%) применяют специализированные онлайн-сервисы записи.
\end{itemize}

Полученные результаты показывают, что в локальной практике доминируют неавтоматизированные и слабо интегрированные способы организации расписания. Это подтверждает актуальность разработки платформы, обеспечивающей централизованное управление слотами, исключение двойного бронирования и унифицированную коммуникацию с клиентом.

С учетом полученных данных можно выделить три практических подхода к ведению записи: ручной, использование готового онлайн-сервиса и разработка собственной платформы. Сравнение в упрощенном виде представлено в таблице \ref{booking:solutions}.

\begin{xltabular}{\textwidth}{|L{3.0cm}|Y|Y|Y|}
\caption{Сравнение подходов к ведению онлайн-записи\label{booking:solutions}}\\ \hline
\centrow Критерий & \centrow Ручной подход (блокнот, звонки) & \centrow Готовый онлайн-сервис & \centrow Разрабатываемая платформа Easyappoint \\ \hline
\endfirsthead
\continuecaption{Продолжение таблицы \ref{booking:solutions}}
\centrow Критерий & \centrow Ручной подход (блокнот, звонки) & \centrow Готовый онлайн-сервис & \centrow Разрабатываемая платформа Easyappoint \\ \hline
\finishhead
Скорость записи клиента & Низкая: требуется согласование по телефону или в чате & Средняя: клиент может записаться сам & Высокая: слот бронируется онлайн по правилам мастера \\ \hline
Учет занятости мастера & Фрагментарный, часто без единого календаря & Обычно есть базовый календарь сервиса & Полный: учитываются расписание и внешний календарь \\ \hline
Уведомления о записи и отмене & Вручную, зависит от дисциплины мастера & Автоматически в рамках возможностей сервиса & Автоматически через события и интеграции \\ \hline
\end{xltabular}

Для уточнения требований были также рассмотрены конкретные готовые решения, ориентированные на автоматизацию записи в сфере услуг: YCLIENTS, DIKIDI и EasyWeek \cite{yclients,dikidi,easyweek}. Эти системы подтверждают востребованность функций онлайн-записи, клиентской базы, уведомлений и календарного управления, но в то же время являются внешними платформами со своей моделью тарификации, пользовательским интерфейсом и ограничениями настройки бизнес-логики.

\begin{xltabular}{\textwidth}{|L{2.5cm}|L{3.0cm}|Y|Y|}
\caption{Сравнение Easyappoint с существующими онлайн-сервисами\label{booking:analogs}}\\ \hline
\centrow Решение & \centrow Основная ориентация & \centrow Сильные стороны & \centrow Ограничения для Easyappoint \\ \hline
\endfirsthead
\caption*{Продолжение таблицы \ref{booking:analogs}}\\ \hline
\centrow Решение & \centrow Основная ориентация & \centrow Сильные стороны & \centrow Ограничения для Easyappoint \\ \hline
\finishhead
YCLIENTS & CRM и онлайн-запись для салонов, клиник и сетевого бизнеса & Клиентская база, электронный журнал, уведомления, аналитика и управление филиалами & Избыточен для одиночного мастера; CRM-модель и закрытая архитектура не дают показать внутренние контракты и outbox-обмен \\ \hline
DIKIDI & Онлайн-запись и каталог исполнителей в сфере услуг & Мобильное приложение, каталог компаний, виджет записи и базовая автоматизация расписания & Зависимость от внешней платформы и ее правил; доменную модель сложно адаптировать под собственные роли, приглашения и интеграции \\ \hline
EasyWeek & Бронирование и управление расписанием для малого бизнеса & Онлайн-бронирование, календарь, напоминания и публичная страница компании & Закрывает типовой сценарий, но не раскрывает проверку конфликтов, доставку событий и контрактный обмен между сервисами \\ \hline
Easyappoint & Собственная платформа для частного мастера и клиентов по приглашению & Контроль доменной модели, ролей, календарной интеграции, Telegram-уведомлений, контрактов и deployment-инфраструктуры & Требует разработки и сопровождения, но позволяет показать полный цикл проектирования в рамках ВКР \\ \hline
\end{xltabular}

С учетом цели работы -- разработка платформы онлайн-бронирования услуг -- наиболее обоснованным является третий подход, предполагающий создание собственного программного решения.

\subsection{Постановка задач на разработку}

На основании проведенного анализа были сформулированы следующие задачи, которые должна решать разрабатываемая платформа:
\begin{enumerate}
\item Обеспечить регистрацию, аутентификацию и авторизацию пользователей с разграничением ролей мастера, клиента и администратора.
\item Реализовать ведение профилей мастера и клиента, а также механизм подключения клиента к мастеру по пригласительной ссылке.
\item Предоставить мастеру средства настройки рабочей доступности, длительности слота и часового пояса.
\item Выполнять расчет доступных временных интервалов с учетом расписания мастера, уже созданных бронирований и занятости во внешнем календаре.
\item Обеспечить создание, просмотр, изменение и отмену бронирований с учетом прав доступа пользователя.
\item Исключить возможность пересечения активных бронирований у одного мастера.
\item Автоматически уведомлять участников процесса о создании, изменении и отмене записи.
\item Обеспечить журналирование ключевых действий и выдачу эксплуатационных метрик для контроля состояния платформы.
\item Предоставить удобный веб-интерфейс для мастера, клиента и администратора.
\end{enumerate}
